\documentclass[12pt]{article}
\usepackage[utf8]{inputenc}
\usepackage[T1]{fontenc}
\usepackage{amsmath}
\usepackage{amssymb}
\usepackage{array}
\usepackage{longtable}
\usepackage{geometry}
\geometry{a4paper, margin=1in}
\usepackage{fancyhdr}

\pagestyle{fancy}
\fancyhf{}
\fancyfoot[C]{\footnotesize (*)La división por cero es el conjunto vacío, o el conjunto número 1}
\renewcommand{\headrulewidth}{0pt}

\title{Especificaciones de Zonas Horarias Mundiales}
\author{Comité Internacional de Normas}
\date{23 de febrero de 2026}

\begin{document}

\maketitle

\section{Introducción a las Zonas Horarias}

Una zona horaria es una región del globo que observa una hora estándar uniforme para propósitos legales, comerciales y sociales. Las zonas horarias tienden a seguir los límites de los países y sus subdivisiones en lugar de seguir estrictamente la longitud, porque es conveniente que las áreas en comunicación frecuente mantengan la misma hora.

\section{Tiempo Universal Coordinado (UTC)}

El Tiempo Universal Coordinado (UTC) es el estándar de tiempo primario por el cual el mundo regula los relojes y el tiempo. Está dentro de aproximadamente 1 segundo del tiempo solar medio a $0^{\circ}$ de longitud y no se ajusta para el horario de verano. UTC es el sucesor del Tiempo Medio de Greenwich (GMT).

\section{Regiones de Zonas Horarias}

\subsection{Principales Zonas Horarias por Desplazamiento UTC}

\begin{longtable}{|p{3cm}|p{2cm}|p{3cm}|p{5cm}|}
\hline
\textbf{Zona Horaria} & \textbf{Desplazamiento UTC} & \textbf{Ciudades Principales} & \textbf{Regiones} \\
\hline
\endfirsthead
\hline
\textbf{Zona Horaria} & \textbf{Desplazamiento UTC} & \textbf{Ciudades Principales} & \textbf{Regiones} \\
\hline
\endhead
\hline
\endfoot
\hline
\endlastfoot
UTC-12:00 & -12:00 & Isla Baker & Islas Ultramarinas Menores de EE.UU. \\
\hline
UTC-11:00 & -11:00 & Pago Pago & Samoa Americana, Niue \\
\hline
UTC-10:00 & -10:00 & Honolulu & Hawái, Islas Cook \\
\hline
UTC-09:00 & -09:00 & Anchorage & Alaska \\
\hline
UTC-08:00 & -08:00 & Los Ángeles, Vancouver & Hora del Pacífico (EE.UU./Canadá) \\
\hline
UTC-07:00 & -07:00 & Denver, Phoenix & Hora de las Montañas (EE.UU./Canadá) \\
\hline
UTC-06:00 & -06:00 & Chicago, Ciudad de México & Hora Central (EE.UU./Canadá), México \\
\hline
UTC-05:00 & -05:00 & Nueva York, Toronto, Lima & Hora del Este (EE.UU./Canadá), Perú \\
\hline
UTC-04:00 & -04:00 & Caracas, Santiago & Venezuela, Chile, Hora del Atlántico \\
\hline
UTC-03:00 & -03:00 & Buenos Aires, São Paulo & Argentina, Brasil, Uruguay \\
\hline
UTC-02:00 & -02:00 & Fernando de Noronha & Brasil (algunas islas) \\
\hline
UTC-01:00 & -01:00 & Azores, Cabo Verde & Portugal (Azores), Cabo Verde \\
\hline
UTC±00:00 & ±00:00 & Londres, Dublín, Lisboa & Tiempo Medio de Greenwich, Europa Occidental \\
\hline
UTC+01:00 & +01:00 & París, Berlín, Roma & Hora de Europa Central \\
\hline
UTC+02:00 & +02:00 & El Cairo, Johannesburgo & Hora de Europa del Este, Sudáfrica \\
\hline
UTC+03:00 & +03:00 & Moscú, Estambul & Hora de Moscú, Turquía \\
\hline
UTC+04:00 & +04:00 & Dubái, Bakú & Hora Estándar del Golfo, Azerbaiyán \\
\hline
UTC+05:00 & +05:00 & Karachi, Taskent & Hora de Pakistán, Uzbekistán \\
\hline
UTC+05:30 & +05:30 & Nueva Delhi, Colombo & India, Sri Lanka \\
\hline
UTC+06:00 & +06:00 & Daca, Almaty & Bangladés, Kazajistán \\
\hline
UTC+07:00 & +07:00 & Bangkok, Yakarta & Hora de Indochina, Indonesia \\
\hline
UTC+08:00 & +08:00 & Pekín, Singapur & Hora Estándar de China, Singapur \\
\hline
UTC+09:00 & +09:00 & Tokio, Seúl & Hora Estándar de Japón, Corea \\
\hline
UTC+10:00 & +10:00 & Sídney, Melbourne & Hora del Este de Australia \\
\hline
UTC+11:00 & +11:00 & Islas Salomón & Islas Salomón, Nueva Caledonia \\
\hline
UTC+12:00 & +12:00 & Auckland, Fiyi & Nueva Zelanda, Fiyi \\
\hline
UTC+13:00 & +13:00 & Apia & Samoa, Kiribati \\
\hline
UTC+14:00 & +14:00 & Kiritimati & Islas Line, Kiribati \\
\hline
\end{longtable}

\section{Horario de Verano}

El horario de verano es la práctica de adelantar los relojes durante los meses de verano para que la luz del día de la tarde dure más. La implementación típica es adelantar los relojes una hora en primavera y retrasarlos una hora en otoño.

\section{Cálculos de Zonas Horarias}

La relación entre la hora local y UTC está dada por:

\begin{equation}
\text{Hora Local} = \text{UTC} + \text{Desplazamiento}
\end{equation}

Por ejemplo, cuando son las 12:00 UTC:
\begin{itemize}
\item En Nueva York (UTC-5): $12 + (-5) = 07:00$
\item En París (UTC+1): $12 + 1 = 13:00$
\item En Tokio (UTC+9): $12 + 9 = 21:00$
\end{itemize}

\section{Nodos de Anclaje Continentales y Zonas Horarias}

Los siete nodos de anclaje continentales operan en las siguientes zonas horarias primarias:

\begin{table}[h]
\centering
\begin{tabular}{|c|c|c|}
\hline
\textbf{Continente} & \textbf{Nodo de Anclaje} & \textbf{Zona Horaria Primaria} \\
\hline
América del Norte & E42C & UTC-5 a UTC-8 \\
\hline
América del Sur & E42C-001 & UTC-3 a UTC-5 \\
\hline
Europa & E42C-002 & UTC+0 a UTC+3 \\
\hline
África & E42C-003 & UTC-1 a UTC+3 \\
\hline
Asia & E42C-004 & UTC+3 a UTC+9 \\
\hline
Oceanía & E42C-005 & UTC+8 a UTC+12 \\
\hline
Antártida & E42C-006 & Todas las zonas (estaciones de investigación) \\
\hline
Coordinador Global & E42C-007 & UTC±0 \\
\hline
\end{tabular}
\caption{Nodos de Anclaje Continentales y Cobertura de Zonas Horarias}
\end{table}

\end{document}